%% ============================================================
%% Section 4: Governance and Regulation
%% Source material: C10 (§4.1) + C11 (§4.2)
%% ============================================================

\section{Governance and Regulation}
\label{sec:governance}

The threats catalogued in Section~\ref{sec:threats} have prompted a global governance response characterized by divergent jurisdictional philosophies, rapid institutional proliferation, and a persistent gap between regulatory frameworks and technological deployment velocity. This section examines government policies, international cooperation, and industry self-regulation. Technical implementation details (e.g., C2PA architecture, watermarking algorithms) are deferred to Section~\ref{sec:countermeasures}.


%% ------------------------------------------------------------
\subsection{Government Policies}
\label{sec:gov-policy}
%% ------------------------------------------------------------

\subsubsection{European Union}
\label{sec:gov-eu}

The EU pursues a prescriptive, rights-centric, rule-based approach anchored in multiple overlapping instruments.

\paragraph{EU AI Act (Regulation 2024/1689).}
The EU AI Act, entered into force August~1, 2024, establishes a risk-tiered framework classifying AI systems into four categories: \emph{unacceptable risk} (banned---e.g., social scoring, subliminal manipulation), \emph{high risk} (strict conformity assessments---e.g., AI in recruitment, critical infrastructure), \emph{limited risk} (transparency obligations---chatbots, deepfake generators), and \emph{minimal risk} (unregulated under the Act)~\cite{euaiact2024}. For AIGC, the most significant provisions reside in the limited-risk tier, which mandates transparency to ensure users recognize interactions with synthetic entities.

Article~50 establishes stringent transparency requirements for synthetic content: providers must ensure AI-generated outputs are marked in machine-readable format and detectable as artificial; deployers must disclose deepfakes even when the underlying content is lawful~\cite{euaiact_art50}. The deepfake labeling obligations become fully enforceable on August~2, 2025, alongside governance rules for General Purpose AI (GPAI) models. Full application of the Act occurs August~2, 2026. Penalties reach up to \euro35~million or 7\% of global annual turnover~\cite{euaiact2024}.

\paragraph{Code of Practice on Transparency.}
The EU AI Office facilitated a Code of Practice on Transparency of AI-Generated Content to operationalize Article~50~\cite{eucode2025}. The first draft (December 2025) mandated multi-layered marking---C2PA metadata, imperceptible watermarks, and visible labels. The streamlined second draft (March 2026) proposed a uniform interactive ``EU Icon'' and free-of-charge detection APIs~\cite{eucode2026draft2}. Adherence serves as formal presumption of conformity with Article~50.

\paragraph{DSA and GDPR.}
The Digital Services Act (DSA) governs AIGC \emph{dissemination}, requiring Very Large Online Platforms ($>$45M EU users) to conduct systemic risk assessments encompassing AI-generated disinformation and deepfakes~\cite{dsa2024}. The GDPR continues to constrain AI training data: CNIL and EDPB guidance (2025--2026) stipulates that if personal data is demonstrably memorized by a model, the right to erasure applies, subject to a balancing test weighing privacy impact against retraining costs~\cite{cnil2025,edpb2024}. Models developed with unlawfully processed data may face destruction orders.


\subsubsection{United States}
\label{sec:gov-us}

The US presents a fragmented landscape characterized by federal executive volatility, voluntary standards, and aggressive state-level gap-filling.

\paragraph{Federal executive volatility.}
Executive Order 14110 (October 2023) imposed safety reporting requirements on frontier AI developers, directed watermarking standards development, and mandated misinformation risk assessment~\cite{eo14110}. This was revoked by Executive Order 14179 (January 2025), which characterized the prior framework as ``onerous government control'' hindering innovation~\cite{eo14179}. The subsequent ``America's AI Action Plan'' (July 2025) shifted federal focus entirely toward deregulation, infrastructure expansion, and ``ideological neutrality'' in LLMs~\cite{aiactionplan2025}.

\paragraph{NIST frameworks.}
NIST provides voluntary technical governance through the AI Risk Management Framework (AI RMF 1.0), expanded via the Generative AI Profile (NIST-AI-600-1, 2024) addressing hallucinations, deepfakes, and provenance~\cite{nistrmf,nistgenai600}. The Cybersecurity Framework Profile for AI (NIST IR 8596, December 2025) integrates AI risks into CSF 2.0 across three pillars: securing AI systems, using AI for detection, and defending against adversarial AI use~\cite{nist8596}.

\paragraph{State-level regulation.}
California's SB~53 (effective January 2026) is the first US law targeting frontier model \emph{developers}, requiring public disclosure of catastrophic risk protocols and mandatory incident reporting~\cite{casb53}. Texas's TRAIGA (effective January 2026) prohibits AIGC for behavioral manipulation, discrimination, and non-consensual deepfakes~\cite{texastraiga}. New York mandates disclosures for political deepfakes. At the federal level, the DEFIANCE Act establishes civil remedies for deepfake NCII victims~\cite{defianceact}, the Take It Down Act (May 2025) criminalizes digital forgeries with 48-hour platform takedown mandates~\cite{takeitdown2025}, and the AI Fraud Deterrence Act doubles penalties for AI-assisted financial fraud~\cite{aifraudact}.


\subsubsection{China}
\label{sec:gov-china}

China operates the most rapid, centrally administered AIGC governance framework, driven by dual mandates of market dominance and sovereign information control.

\paragraph{Phased regulatory architecture.}
The Deep Synthesis Management Provisions (effective January 2023) require real-name identity verification, content labeling, and content management systems for all deep-synthesis services~\cite{chinadeepsynthesis}. The Interim Measures for the Management of Generative AI Services (effective August 2023) mandate that public-facing AIGC services ensure training data legality, actively monitor outputs, and uphold ``core socialist values''~\cite{chinainterim2023}. Providers bear legal responsibility for model outputs.

\paragraph{Algorithm filing system.}
Under the Algorithmic Recommendation Management Provisions, any AI service with ``public opinion attributes or social mobilization capabilities'' must conduct internal risk assessment and file its algorithm with the Cyberspace Administration of China (CAC) before public deployment~\cite{chinaalgorithm}. This operates as a mandatory ex-ante licensing regime, granting the state complete visibility into the commercial AI landscape.

\paragraph{TC260 security requirements.}
TC260's ``Basic Security Requirements for Generative AI Services'' (TC260-003, 2024) define 31 specific security risks across five categories, with particular emphasis on ``corpus safety'' (training data integrity) and political compliance~\cite{tc260}. Developers must statistically demonstrate that generated samples do not contain enumerated risks to pass CAC filing. Active enforcement includes the ``Sword Pujiang'' action (late 2025) that delisted 54 non-compliant applications~\cite{chinaenforcement2025}.


\subsubsection{International Cooperation}
\label{sec:gov-intl}

\paragraph{UK AI Security Institute.}
The UK AI Safety Institute, rebranded as the AI Security Institute (AISI) in early 2025, conducts empirical evaluations of frontier models. Its December 2025 ``Frontier AI Trends Report'' found AI performance in sensitive domains doubling every eight months, with frontier models surpassing PhD-level baselines in chemistry and biology by up to 60\%~\cite{aisi2025trends}.

\paragraph{International AI Safety Reports.}
The International AI Safety Reports (2025, 2026), chaired by Yoshua Bengio with 100+ experts from 30 nations, provide authoritative scientific consensus~\cite{aisafetyreport2025,aisafetyreport2026}. The 2026 report documented a 92.7\% year-over-year increase in AI-assisted data exfiltration. The US declined to endorse the 2026 edition~\cite{time2026usretreat}.

\paragraph{Summit trajectory.}
The AI Safety Summit series---Bletchley Park (November 2023, frontier safety focus), Seoul (May 2024, innovation and inclusivity), Paris (February 2025, economic adoption and sustainability), New Delhi (February 2026, Global South access and democratization)---reveals a fracturing of priorities between safety-focused and innovation-focused camps, with the US and UK declining to sign the Paris declaration~\cite{bletchley2023,seoul2024,paris2025}.

\paragraph{G7 and OECD.}
The G7 Hiroshima AI Process produced the International Code of Conduct for Advanced AI Systems, urging pre-deployment risk assessments, content authentication, and IP respect~\cite{hiroshima2023}. The OECD AI Principles, updated May 2024 to address generative AI, elevated ``information integrity'' as a core concern and are now adhered to by 47 governments~\cite{oecdai2024}.

\paragraph{Comparative analysis.}
The global governance landscape has fractured into three paradigms: the \textbf{EU model} (prescriptive, rights-centric, rule-based---strong protection but high compliance friction), the \textbf{US model} (voluntary, market-driven, principle-based---innovation-facilitating but creating regulatory vacuum and state-level patchwork), and the \textbf{China model} (centrally administered, application-specific, ideology-entwined---agile but subordinating model capability to political control). International efforts represent attempts to bridge these irreconcilable philosophies.


%% ------------------------------------------------------------
\subsection{Industry Practices}
\label{sec:gov-industry}
%% ------------------------------------------------------------

\subsubsection{Model Provider Safety Frameworks}
\label{sec:gov-frameworks}

The Seoul AI Safety Summit (May 2024) catalyzed a concentrated burst of self-regulation: twelve frontier AI companies published safety frameworks by December 2025~\cite{metr2025frameworks}. All share a common architecture of capability thresholds triggering escalating safeguards, but diverge significantly in specificity and governance.

Anthropic's Responsible Scaling Policy (RSP v3.0, February 2026) defines AI Safety Levels (ASL-1 through ASL-4+) modeled on biosafety levels, with ASL-3 activated for CBRN-related risks requiring robust jailbreak resistance~\cite{anthropicrsp3}. OpenAI's Preparedness Framework v2.0 (April 2025) streamlines to two thresholds (``High'' and ``Critical'') but controversially removed persuasion risk tracking~\cite{openaiprepv2}. Google DeepMind's FSF v3.0 (September 2025) introduced Critical Capability Levels including explicit misalignment evaluation~\cite{googlefsf3}. Meta's framework takes an ``outcomes-led approach,'' evaluating whether models \emph{uniquely enable} catastrophic scenarios~\cite{metafrontier2025}. All frameworks remain \textbf{voluntary, self-enforced, and subject to self-assessed thresholds} with no independent verification.

\paragraph{FLI AI Safety Index.}
The Future of Life Institute's AI Safety Index (Winter 2025) provides the most systematic external assessment~\cite{fli2025winter}. Anthropic ranked first (C+, 2.67), followed by OpenAI (C+, 2.31) and Google DeepMind (C, 2.08); xAI, Meta, DeepSeek, and Alibaba Cloud received D or D- grades. Most strikingly, \textbf{no company scored above D in Existential Safety}---all are ``racing toward AGI without presenting any explicit plans for controlling such technology''~\cite{fli2025winter}.

\paragraph{Red teaming.}
Red teaming has matured into three modalities: \emph{human expert red teaming} (Anthropic's 38,961-attack dataset~\cite{ganguli2022redteam}; partnerships with NNSA and national AI Safety Institutes), \emph{automated tools} (Microsoft PyRIT~\cite{pyrit2024}, NVIDIA Garak~\cite{garak2024}), and \emph{external programs} (OpenAI bug bounty at \$100K maximum; DEF~CON 31 challenge with 2,244 hackers evaluating eight LLMs~\cite{defcon2023generativeredteam}).


\subsubsection{Platform Governance and Standards}
\label{sec:gov-platform}

Major platforms have converged on \emph{labeling rather than removal} as the primary governance mechanism for AI-generated content.

\textbf{Meta} evolved through ``Imagined with AI'' $\to$ ``Made with AI'' $\to$ ``AI Info'' labels (2023--2024), using C2PA and IPTC metadata for third-party content detection, while prohibiting generative AI ad tools for political ads~\cite{metalabeling2024}. \textbf{Google/YouTube} deploys SynthID watermarking across images, audio, text, and video, with over 10 billion pieces of content watermarked by May 2025; YouTube requires creators to disclose synthetic content~\cite{synthid2025}. \textbf{TikTok} is the most restrictive: automatic C2PA detection, complete ban on AI political ads, and removal of 51,000 synthetic media videos while labeling over 1.3 billion~\cite{tiktok2023labels}. \textbf{X} represents the opposite extreme with no significant AI-specific policy updates since 2020 and fewer than 9\% of Community Notes reaching the general audience~\cite{xpolicy2020}.

\paragraph{C2PA and Content Credentials.}
The Coalition for Content Provenance and Authenticity (C2PA), now at specification v2.2 (May 2025), has achieved broad adoption: hardware embedding (Leica M11-P, Google Pixel~10), AI tool integration (Adobe Firefly, DALL-E~3, Sora), and platform detection (Meta, Google, TikTok, LinkedIn)~\cite{c2pa2025}. The Content Authenticity Initiative has grown to over 5,000 members. CISA endorsed Content Credentials in January 2025. Despite momentum, metadata stripping during sharing and uneven ecosystem adoption remain fundamental challenges.

\paragraph{Industry alliances.}
The Partnership on AI (126 partners, 16 countries)~\cite{pai}, the Frontier Model Forum (\$10M+ AI Safety Fund)~\cite{fmf2023}, and the White House Voluntary AI Commitments (15 companies, July--September 2023)~\cite{whitehouse2023commitments} provide institutional scaffolding. The Seoul Declaration (10 countries plus EU) and Frontier AI Safety Commitments (16 companies) catalyzed the International Network of AI Safety Institutes. However, a one-year assessment found companies ``nowhere near where we need them to be in terms of good governance''~\cite{mitreview2024assessment}.


\subsubsection{Agentic AI Governance Gap}
\label{sec:gov-agentic}

The governance gap for agentic AI is structural, not merely temporal. Current frontier safety frameworks, the NIST AI RMF, ISO/IEC 42001, and the EU AI Act contain \textbf{zero references to ``agent,'' ``agentic,'' or autonomous AI systems}~\cite{zenity2025blindspot}. This is a critical mismatch: AI agents introduce qualitatively different risks---autonomous real-world action, dynamic tool access, multi-step planning with cascading failures, persistent memory vulnerable to manipulation, and multi-agent delegation fragmenting accountability.

The deployment timeline is stark. Anthropic launched Computer Use (October 2024), OpenAI's Operator followed (January 2025), Claude Code reached GA (May 2025, \$1B annualized revenue by November), and Codex accumulated 1.6M weekly users by March 2026~\cite{anthropic2024computeruse,openaicopilot}. \textbf{None launched under any agent-specific governance standard.} Consequences materialized when Anthropic disclosed that a Chinese state-sponsored actor used Claude Code for what it described as the first AI-orchestrated cyber espionage campaign, with AI executing 80--90\% of operations autonomously~\cite{anthropic2025gtg1002}.

MCP exemplifies the ``deploy capability first, retrofit security later'' pattern. The original specification used ``SHOULD'' rather than ``MUST'' for security controls, and over 13,000 servers launched in 2025 alone, many without authentication~\cite{zenity2025mcp}. Security iteratively improved---OAuth 2.1 (March 2025), Resource Server classification (June 2025), enhanced authorization (November 2025)---but significant gaps persist for autonomous agent scenarios~\cite{aws2025mcpauth}.

Emerging standards are chasing a moving target. The \textbf{OWASP Top~10 for Agentic Applications} (December 2025) codifies ten agent-specific risks~\cite{owasp2025agentic}. \textbf{Singapore's IMDA} published the world's first government agentic AI governance framework (January 2026)~\cite{imda2026agentic}. \textbf{NIST} launched its AI Agent Standards Initiative (February 2026)~\cite{nist2026agentstandards}. The Cloud Security Alliance published an Agentic Trust Framework applying Zero Trust to AI agents~\cite{csa2026agentic}. Yet standardization remains 16+ months behind deployment: the first agent products shipped October 2024, the first government frameworks arrived January 2026.

\begin{table}[htbp]
\centering
\caption{Cross-jurisdiction AIGC governance comparison.}
\label{tab:governance}
\small
\begin{tabular}{p{1.8cm}p{3.2cm}p{2.8cm}p{2.0cm}p{2.8cm}}
\toprule
\textbf{Jurisdiction} & \textbf{Key Legislation} & \textbf{Core Mechanism} & \textbf{Stage} & \textbf{Maximum Penalty} \\
\midrule
EU &
AI Act (2024/1689); DSA; GDPR; Code of Practice &
Risk-tiered regulation; mandatory deepfake labeling; ex-ante conformity &
Phased (full by Aug 2026) &
\euro35M or 7\% global turnover \\
\addlinespace
US &
EO 14179; NIST AI RMF; SB~53 (CA); TRAIGA (TX); Take It Down Act &
Voluntary federal standards; state patchwork; targeted federal bills &
Fragmented &
Varies by state/bill \\
\addlinespace
China &
Deep Synthesis Provisions; Interim Measures; TC260-003 &
Algorithm filing; corpus safety; real-name verification; political compliance &
Fully operational &
Service suspension; delisting \\
\addlinespace
International &
AI Safety Reports; Hiroshima Code; OECD Principles; Summit declarations &
Soft law; voluntary commitments; scientific consensus &
Evolving &
Non-binding \\
\bottomrule
\end{tabular}
\end{table}
